\documentclass{beamer}

% ---------- Thème et couleurs ----------
\usetheme{Madrid}
\usecolortheme{beaver}
\setbeamertemplate{navigation symbols}{}
\setbeamertemplate{footline}[frame number]

% ---------- Packages ----------
\usepackage[utf8]{inputenc}
\usepackage[T1]{fontenc}
\usepackage[french]{babel}
\usepackage{graphicx}
\usepackage{array}
\usepackage{booktabs}
\usepackage{xcolor}

% ---------- Titre ----------
\title{Combinaison approach Distribution CNN \& MARL }
\author{Supervisé par : AIT OMAR Driss\\
            Préparé par : AIT AHMED OULHAJE Soulaimane}
\institute{Master AIDC — FST Béni Mellal}

% ==========================================================
\begin{document}

% ---------- Page de titre ----------
\begin{frame}
  \titlepage
\end{frame}
% ---------- Slide : Evolution des réseaux IoT ----------
\begin{frame}{Introduction}
\begin{itemize}
    \item L’\textbf{Internet des Objets (IoT)} et l’\textbf’intelligence artificielle (IA) ont profondément transformé notre manière de collecter, traiter et analyser les données.
    \item Ces données sont souvent exploitées par des modèles profonds tels que les \textbf{réseaux de neurones convolutionnels (CNN)} pour des tâches d’analyse avancées :
    \begin{itemize}
        \item vision par ordinateur (surveillance, reconnaissance faciale),
        \item analyse médicale (IRM, radiologie),
        \item systèmes intelligents (véhicules autonomes, smart cities).
    \end{itemize}
\end{itemize}
\end{frame}

\begin{frame}{Problématique}
\begin{itemize}
    \item Traditionnellement, les calculs d’inférence sont délégués à des \textbf{serveurs cloud distants} capables de gérer la lourde charge computationnelle.
    \item Cependant, cette centralisation pose deux limites majeures :
\end{itemize}

\end{frame}

% ---------- Slide : Problème de performance ----------
\begin{frame}{Impact sur les performances}

\textbf{La centralisation de l'inférence CNN entraîne une latence élevée :}

\begin{itemize}
    \item Les données brutes (images, vidéos, signaux) doivent être envoyées au cloud.
    \item Le transfert réseau ajoute un délai important.
    \item Cette latence rend le cloud inadapté aux applications temps réel :
    \begin{itemize}
        \item véhicules autonomes,
        \item vidéosurveillance en direct,
        \item chirurgie robotique à distance,
        \item systèmes de sécurité critiques.
    \end{itemize}
\end{itemize}

\begin{block}{Conclusion}
�� La latence du cloud compromet la performance des services IoT sensibles au temps.
\end{block}

\end{frame}

% ---------- Slide : Problème de confidentialité ----------
\begin{frame}{Impact sur la confidentialité}

\textbf{La centralisation expose fortement les données sensibles :}

\begin{itemize}
    \item Les données brutes quittent leur source pour être traitées dans des serveurs distants (Cloud).
    \item Elles peuvent être interceptées ou surveillées pendant la transmission.
    \item Un serveur central devient un \textbf{point unique de vulnérabilité}.
\end{itemize}

\begin{block}{Conclusion}
\begin{itemize}
    \item La centralisation augmente fortement le risque de fuite ou de compromission des données.
\end{itemize}
\end{block}

\end{frame}

% ---------- Slide : Pourquoi le cloud ne suffit pas ----------
\begin{frame}{}
\begin{block}{Conclusion}
\begin{itemize}
    \item Le cloud est puissant, mais trop loin et trop risqué pour les applications IoT critiques. Face à ces défis, l’\textbf{Edge Computing} et l’\textbf{apprentissage distribué} émergent comme solutions pour rapprocher l’IA des sources de données.
\end{itemize}
\end{block}
\end{frame}

% ---------- Slide : Définition du Edge Computing ----------
\begin{frame}{Qu'est-ce que le Edge Computing ?}

\begin{itemize}
    \item Le \textbf{Edge Computing} est un paradigme où le traitement des données 
    est réalisé \textbf{au plus près de leur source} (caméras, capteurs, drones, smartphones, etc.).
    \item Au lieu d'envoyer toutes les données vers le cloud, 
    une partie de l'inférence et du calcul est effectuée localement.
    \item Objectifs :
    \begin{itemize}
        \item réduire la latence ;
        \item diminuer le trafic réseau ;
        \item améliorer la confidentialité en gardant les données sensibles sur site.
    \end{itemize}
\end{itemize}

\end{frame}

\begin{frame}{Limites du Edge Computing}
\begin{itemize}
    \item Moins de latence mais capacité de calcul insuffisante.
    \item Incapacité à héberger des CNN profonds.
    \item Contraintes : énergie, mémoire, bande passante.
\end{itemize}
\end{frame}

% ---------- Slide 1 ----------
\begin{frame}{Pourquoi l’apprentissage distribué est une solution ?}

Le \textbf{l’apprentissage distribué } consiste à faire collaborer plusieurs appareils IoT 
pour exécuter ensemble une tâche trop lourde pour un seul dispositif.

\vspace{0.3cm}
Cette approche permet de dépasser les limites du cloud et du edge computing 
dans les environnements IoT modernes.

\end{frame}

% ---------- Slide : Privacy-Aware Distributed CNN ----------
\begin{frame}{Pourquoi l’apprentissage distribué est une solution ?}

\textbf{Principe :}
\begin{itemize}
    \item Distribuer les couches ou les feature maps et Effectuer le calcul d'un CNN entre plusieurs appareils IoT près de la source.
    \item Fragmenter les données afin que chaque nœud ne voie qu'une petite partie de l'information.
    \item Brouiller ou disperser les représentations internes pour empêcher la reconstruction des données sensibles.
\end{itemize}

\begin{block}{Objectif}
Préserver la confidentialité tout en permettant l'exécution distribuée d'un CNN dans un réseau IoT.
\end{block}

\end{frame}

% ---------- Slide 1 : Qu'est-ce que le MARL ? ----------
\begin{frame}{Qu'est-ce que le MARL ?}

\textbf{MARL : Multi-Agent Reinforcement Learning}

\begin{itemize}
    \item Extension du RL où plusieurs agents apprennent simultanément.
    \item Chaque agent observe l'environnement, prend une décision
    et reçoit une récompense.
    \item Les agents apprennent une stratégie optimale en tenant compte
    les uns des autres.
\end{itemize}

\begin{block}{Caractéristique clé}
�� Le comportement optimal dépend des décisions des autres agents.
\end{block}

\end{frame}

% ---------- Slide 2 : Pourquoi le MARL est nécessaire ? ----------
\begin{frame}{Pourquoi utiliser le MARL dans les réseaux IoT ?}

Le problème présente 4 caractéristiques clés :

\begin{enumerate}
    \item Plusieurs entités indépendantes (hôpital, aéroport, banque...).
    \item Objectifs différents mais ressources IoT partagées.
    \item Environnement dynamique et non déterministe.
    \item Aucune autorité centrale pour coordonner les décisions.
\end{enumerate}

\begin{block}{Conclusion}
�� Ces conditions correspondent parfaitement aux scénarios MARL.
\end{block}

\end{frame}

\begin{frame}{Différence entre RL mono-agent et MARL}

\centering
\renewcommand{\arraystretch}{1.5}

\begin{tabular}{|l|p{3.9cm}|p{3.9cm}|}
\hline
\textbf{Critère} & \textbf{RL mono-agent} & \textbf{MARL} \\
\hline
Architecture & Centralisée & Distribuée \\
\hline
Scalabilité & Peu scalable & Hautement scalable \\
\hline
Prise de décision & Décisions globales lentes & Décisions locales rapides \\
\hline
Robustesse & Fragile (point unique de défaillance) & Robuste \\
\hline
Complexité & Simple à implémenter & Réaliste industriellement \\
\hline
\end{tabular}

\vspace{0.4cm}

\textbf{Conclusion :}  
Le MARL est plus adapté aux environnements IoT distribués et dynamiques.

\end{frame}

\begin{frame}{Formulation mathématique de $L_{IoT}$}

\[
\min_{A_{r,l,p}^i} \; L_{IoT}
\]

avec :

\[
L_{IoT} =
\sum_{r \in RQ_n}
\left(
\sum_{l=2}^{L_r-2}
\max_{i,j}
\left(
\frac{O^{r,l-1}_{i,j}}{\rho_i}
+
t^{r,l,j}_c
\right)
+ t_s + t_f
\right)
\]

\begin{itemize}
  \item $L_{IoT}$ n’est pas un simple temps de calcul
  \item C’est une \textbf{fonction de coût globale}
  \item Elle permet de :
  \begin{itemize}
    \item distribuer les couches CNN sur les devices IoT
    \item minimiser le temps d’inférence bout-en-bout
    \item respecter les contraintes de ressources et de confidentialité
  \end{itemize}
\end{itemize}
\end{frame}

\begin{frame}{Conclusion}
\footnotesize
\textbf{Contraintes de ressources} \\
\medskip
\textbf{Contrainte de communication} \\
\medskip
\textbf{Contraintes structurelles du CNN} \\
\medskip
\textbf{Contraintes temporelles et confidentialité}

\begin{block}{Message clé}
$L_{IoT}$ formalise le compromis fondamental entre calcul, communication et contraintes système dans l’inférence CNN distribuée.
\end{block}

\bigskip
\centering
\textbf{Ce qui justifie son optimisation par RL et MARL.}
\end{frame}

\begin{frame}{De la formulation centralisée au MARL}

\small
\textbf{L\'objectif Global :}
\begin{itemize}
  \item Tous les agents prennent des décisions
  \item Mais partagent les mêmes contraintes globales
  \item $\Rightarrow$ conflits directs sur les ressources
  \item $\Rightarrow$ optimisation centralisée peu scalable
\end{itemize}

\medskip
\textbf{Avec décomposition duale :}
\begin{itemize}
  \item Chaque agent résout son propre problème local
  \item Les contraintes globales sont relaxées dans l’objectif
  \item Introduction de multiplicateurs duaux (prix)
\end{itemize}

\medskip
\begin{block}{Rôle des multiplicateurs}
\begin{itemize}
  \item pénalisent l’utilisation excessive des ressources
  \item reflètent la congestion du système
  \item forcent une convergence globale
\end{itemize}
\end{block}

\end{frame}

\begin{frame}{Décomposition de l’objectif global}

\footnotesize
Le problème d’optimisation peut être interprété comme :

\[
\min_{A} \;
\underbrace{
\sum_{r\in RQ_n}
\sum_{l=2}^{L_r-2}
\max\!\left(
\frac{O^{r,l-1}_{i,j}}{\rho_i}
+
t^{r,l,j}_c
\right)
+ t_s + t_f
}_{\text{(A) Latence bout-en-bout}}
\]

\[
+
\underbrace{
\sum_{r,l,p}\sum_{i\in\mathcal{I}}
A^i_{r,l,p}(m^r_{l,p}+c^r_{l,p})
-
(\bar m_i+\bar b_i+\bar c_i)
}_{\text{(B) Pénalisation des ressources (relaxation duale)}}
\]

\[
+
\underbrace{
\sum_{a\in N,\,a\neq n}
\sum_{r^\*\in RQ_a}
\sum_{l^\*,p^\*}
g^{\,r,l,p}_{i,r^\*,l^\*,p^\*}
}_{\text{(C) Couplage inter-agents (coordination MARL)}}
\]

\medskip
\begin{block}{Lecture clé}
\begin{itemize}
  \item (A) : objectif primal non décomposable (chemin critique)
  \item (B) : contraintes intégrées comme coûts (dualité)
  \item (C) : interactions entre agents (coordination MARL)
\end{itemize}
\end{block}

\end{frame}

\begin{frame}{(A) Latence bout-en-bout : objectif primal}

\small
\[
\sum_{r\in RQ_n}
\left(
\sum_{l=2}^{L_r-2}
\max\!\left(
\frac{O^{r,l-1}_{i,j}}{\rho_i}
+
t^{r,l,j}_c
\right)
+ t_s + t_f
\right)
\]

\begin{block}{Rôle}
\begin{itemize}
  \item Capture le \textbf{chemin critique} du CNN distribué
  \item Intègre calcul et communication
  \item Reflète la réalité physique des systèmes distribués
\end{itemize}
\end{block}

\medskip
\textbf{Terme non décomposable :} il impose une coordination globale.
\end{frame}

\begin{frame}{(B) Pénalisation des ressources : relaxation duale}

\small
\[
\sum_{r,l,p}\sum_{i\in\mathcal{I}}
A^i_{r,l,p}
\big(
m^r_{l,p}+c^r_{l,p}
\big)
-
(\bar m_i+\bar c_i+\bar b_i)
\]

\begin{block}{Rôle}
\begin{itemize}
  \item Intègre les contraintes dans l’objectif
  \item Transforme les contraintes globales en \textbf{coûts locaux}
  \item Joue le rôle de multiplicateurs duaux (prix)
\end{itemize}
\end{block}

\medskip
\textbf{Effet :} dépassement de ressources $\Rightarrow$ pénalité.
\end{frame}

\begin{frame}{(C) Couplage inter-agents : coordination MARL}

\small
\[
\sum_{a \in \mathcal{N},\, a \neq n}
\sum_{r^{*} \in RQ_a}
\sum_{l^{*},\,p^{*}}
g^{\,r,l,p}_{\,i,\,r^{*},\,l^{*},\,p^{*}}
\]


\begin{block}{Rôle}
\begin{itemize}
  \item Modélise les conflits entre agents
  \item Capture la compétition pour les ressources partagées
  \item Assure une coordination implicite entre agents
\end{itemize}
\end{block}

\medskip
Sans ce terme → agents égoïstes \\

Avec ce terme → agents coordonnés \\ 

\textbf{Ce terme justifie l’utilisation du MARL.} \\
\end{frame}

\begin{frame}{MARL pour l’inférence CNN distribuée}

\small
Afin de faire face au caractère \textbf{dynamique et distribué} du système IoT,
le problème de distribution du CNN est modélisé comme un
\textbf{processus de décision markovien partiellement observable multi-agent (POMDP)}.

\medskip
\begin{itemize}
  \item Chaque agent représente une \textbf{source de données indépendante}
  \item Les agents disposent :
  \begin{itemize}
    \item d’observations partielles de l’environnement global
    \item d’informations locales privées liées à la confidentialité
  \end{itemize}
  \item Les décisions sont prises de manière \textbf{locale et autonome}
\end{itemize}

\medskip
\begin{block}{Objectif}
Le système MARL permet une \textbf{adaptation en ligne} aux changements du réseau
tout en minimisant la \textbf{latence d’inférence} et le \textbf{surcoût induit par la confidentialité}.
\end{block}

\end{frame}


\end{document} 